\documentclass[tikz, border=10pt]{standalone}
\usepackage{graphicx}
\usepackage{tikz}
\usetikzlibrary{calc,arrows.meta}
\usepackage{pgfplotstable}

\title{Thesis Plots}
\author{Mojtaba Kanani}
\date{December 2025}

% --------------------------------------------------
% Histogram style switches
% --------------------------------------------------
\newif\ifHistogramFilled        % boolean flag: filled vs outline

\begin{document}

% --------------------------------------------------
% Configuration
% --------------------------------------------------
\def\TotalFrames{100}              % total number of frames you have
\def\FrameRanges{40/60}            % ranges to SHOW: a/b, c/d, e/e for single frames

% Row bases (in tikz y-units)
\def\FrameNumberBase{0}            % y position of frame-number row
\def\HistogramBase{1.1}            % y position of histogram baseline

% Histogram configuration
\def\HistogramScale{4}             % vertical scale factor for bar height
\def\HistogramColor{black}         % color of bars (draw and/or fill)
\def\HistogramBarFraction{0.95}    % fraction of image width used for bar (0<...<=1)
\HistogramFilledtrue               % set to \HistogramFilledfalse for outline bars

% Highlight windows configuration
% Windows in frame indices, e.g. 45/50, 48/52 etc.
\def\HighlightRanges{41/45, 47/51, 53/57} % example: two windows
\def\HighlightColor{cyan}
\def\HighlightLineWidth{0.6pt}
\def\HighlightFillOpacity{0.1}     % 0 = transparent (outline only)
\def\HighlightBottom{-0.9}         % vertical span of highlight
\def\HighlightTop{3}

% Horizontal measurement lines configuration
% Each entry: start/stop frame range
\def\MeasurementRanges{41/45, 41/46}
% Corresponding colors (must match number of ranges)
\def\MeasurementColors{cyan, magenta}
% Corresponding labels (must match number of ranges)
\def\MeasurementLabels{Window Size = 5, Stride = 6}
% Vertical positions (negative = below frame numbers)
\def\MeasurementLevels{-1.5, -2.5}

% CSV with one column: row k -> frame (k+1) or whatever convention you use
\pgfplotstableread[col sep=comma]{Data/eventfulness_1-1_15fps.csv}\FrameValues

\begin{tikzpicture}[x=0.5cm,y=0.5cm]

  % -----------------------------
  % Row labels (use the bases)
  % -----------------------------
  \node[anchor=east, font=\scriptsize] at (-0.3, \FrameNumberBase - 0.3) {Frame Number};
  \node[anchor=east, font=\scriptsize] at (-0.3, \HistogramBase + 0.1) {Eventfulness};

  % state for compressed x-position and last shown frame index
  \newcount\displayindex
  \newcount\prevEnd
  \displayindex=0
  \prevEnd=0

  % -----------------------------
  % Draw frames + histogram
  % -----------------------------
  \foreach \start/\stop in \FrameRanges {%

    % gap before this range? draw "..."
    \ifnum\start>\numexpr\prevEnd+1\relax
      \node[anchor=south, inner sep=0pt, font=\scriptsize]
        at (\the\displayindex+0.5, 0.4) {$\cdots$};
      \global\advance\displayindex by 1
    \fi

    % frames in [start, stop]
    \foreach \i in {\start,...,\stop} {%

      % filenames: frame_00040.png etc.
      \ifnum\i<10
        \edef\fname{frame_0000\i.png}%
      \else
        \ifnum\i<100
          \edef\fname{frame_000\i.png}%
        \else
          \edef\fname{frame_00\i.png}%
        \fi
      \fi

      % image node at compressed position
      \node[anchor=south west,inner sep=0pt] (img\i)
        at (\the\displayindex, 0)
        {\includegraphics[height=0.5cm]{Pictures/Frames/\fname}};

      % frame number row: centered under the frame slot
      \node[anchor=north, inner sep=0.1cm, font=\scriptsize]
        at (\the\displayindex+0.5, \FrameNumberBase) {\i};

      % histogram bar for this frame
      % row index in table is i (adjust if your CSV is 0-based or has headers)
      \pgfplotstablegetelem{\numexpr\i\relax}{0}\of{\FrameValues}
      \edef\thisvalue{\pgfplotsretval}%
      \pgfmathsetmacro{\barheight}{\HistogramScale * \thisvalue}%

      % bar width as fraction of image width, centered under the image
      \pgfmathsetmacro{\barLeftFrac}{(1 - \HistogramBarFraction)/2}
      \pgfmathsetmacro{\barRightFrac}{(1 + \HistogramBarFraction)/2}

      \coordinate (barL) at ($(img\i.south west)!\barLeftFrac!(img\i.south east)$);
      \coordinate (barR) at ($(img\i.south west)!\barRightFrac!(img\i.south east)$);

      \ifHistogramFilled
        \draw[draw=\HistogramColor, fill=\HistogramColor]
          ($(barL |- 0,\HistogramBase)$)
          rectangle ($(barR |- 0,\HistogramBase+\barheight)$);
      \else
        \draw[draw=\HistogramColor]
          ($(barL |- 0,\HistogramBase)$)
          rectangle ($(barR |- 0,\HistogramBase+\barheight)$);
      \fi

      \global\advance\displayindex by 1
    }%

    \global\prevEnd=\stop\relax
  }%

  % trailing gap after last range? final "..."
  \ifnum\TotalFrames>\numexpr\prevEnd\relax
    \node[anchor=south, inner sep=0pt, font=\scriptsize]
      at (\the\displayindex+0.5, 0.4) {$\cdots$};
  \fi

  % -----------------------------
  % Draw highlight windows + stats
  % -----------------------------
  \makeatletter
  \foreach \hstart/\hstop in \HighlightRanges {%
    % Check that both img<hstart> and img<hstop> exist as TikZ nodes.
    \pgfutil@ifundefined{pgf@sh@ns@img\hstart}{
      % start node missing -> do nothing
    }{
      \pgfutil@ifundefined{pgf@sh@ns@img\hstop}{
        % stop node missing -> do nothing
      }{
        % ===========================
        % Compute avg and max in window
        % ===========================
        \xdef\winSum{0}
        \xdef\winMax{0}
        \newcount\winCount
        \winCount=0

        \foreach \f in {\hstart,...,\hstop}{%
          % Use the SAME indexing as the histogram bars
          \pgfplotstablegetelem{\numexpr\f\relax}{0}\of{\FrameValues}
          \let\thisvalue\pgfplotsretval
          
          % Accumulate sum and find max
          \pgfmathparse{\winSum + \thisvalue}
          \xdef\winSum{\pgfmathresult}
          
          \pgfmathparse{max(\winMax, \thisvalue)}
          \xdef\winMax{\pgfmathresult}
          
          \global\advance\winCount by 1
        }%
        
        % Calculate average
        \pgfmathparse{\winSum / \the\winCount}
        \xdef\winAvg{\pgfmathresult}

        % ===========================
        % Geometry: rectangle and label
        % ===========================
        \coordinate (winL) at (img\hstart.south west);
        \coordinate (winR) at (img\hstop.south east);
        \coordinate (winC) at ($(winL)!0.5!(winR)$);

        % Highlight rectangle
        \draw[
          rounded corners,
          draw=\HighlightColor,
          line width=\HighlightLineWidth,
          fill=\HighlightColor,
          fill opacity=\HighlightFillOpacity
        ]
        ($(winL |- 0,\HighlightBottom)$)
        rectangle
        ($(winR |- 0,\HighlightTop)$);

        % Stats label at top center of window - TWO LINES
        \node[anchor=south, font=\scriptsize, text=\HistogramColor, align=center]
          at ($(winC |- 0,\HighlightTop)$)
          {avg=\pgfmathprintnumber[fixed,precision=2]{\winAvg}\\
           max=\pgfmathprintnumber[fixed,precision=2]{\winMax}};
      }
    }
  }%
  \makeatother

  % -----------------------------
  % Draw measurement lines
  % -----------------------------
  \makeatletter
  % Parse lists into parallel iteration
  \newcount\measurementIndex
  \measurementIndex=0
  
  \foreach \mstart/\mstop in \MeasurementRanges {%
    % Get corresponding color, label, and level
    \pgfmathsetmacro{\currentIndex}{\the\measurementIndex}
    
    % Extract color
    \foreach \col [count=\c from 0] in \MeasurementColors {%
      \ifnum\c=\currentIndex\relax
        \xdef\currentColor{\col}
      \fi
    }%
    
    % Extract label
    \foreach \lbl [count=\l from 0] in \MeasurementLabels {%
      \ifnum\l=\currentIndex\relax
        \xdef\currentLabel{\lbl}
      \fi
    }%
    
    % Extract level
    \foreach \lvl [count=\v from 0] in \MeasurementLevels {%
      \ifnum\v=\currentIndex\relax
        \xdef\currentLevel{\lvl}
      \fi
    }%
    
    % Check that both nodes exist
    \pgfutil@ifundefined{pgf@sh@ns@img\mstart}{}{%
      \pgfutil@ifundefined{pgf@sh@ns@img\mstop}{}{%
        % Get left and right coordinates
        \coordinate (measL) at (img\mstart.south west);
        \coordinate (measR) at (img\mstop.south east);
        \coordinate (measC) at ($(measL)!0.5!(measR)$);
        
        % Draw line with arrows
        \draw[\currentColor, line width=0.8pt, {Rays[n=2]}-{Rays[n=2]}]
          ($(measL |- 0,\currentLevel)$) -- ($(measR |- 0,\currentLevel)$);
        
        % Add label below the line
        \node[anchor=north, font=\scriptsize, text=\currentColor]
          at ($(measC |- 0,\currentLevel)$)
          {\currentLabel};
      }%
    }%
    
    \global\advance\measurementIndex by 1
  }%
  \makeatother

\end{tikzpicture}

\end{document}